% !TEX TS-program = lualatex
\documentclass[11pt]{beamer}
\usepackage{luatexja} % 文字化け防止のためLuaLaTeXを使用
\usepackage{amsmath, amssymb}

% テーマとカラー設定
\usetheme{Madrid}
\usecolortheme{whale}

% スライド番号の設定
\setbeamertemplate{footline}[frame number]

% タイトル情報
\title{テンソル積と自己準同型環による\\有限次Galois拡大の近近代 的・構造論的考察}
\author{}
\date{}

\begin{document}

\frame{\titlepage}

% --- 目次スライド ---
\begin{frame}{目次}
    \tableofcontents
\end{frame}
% -------------------------

\section{概要と解説の目的}
\begin{frame}[allowframebreaks]{概要と解説の目的}
    本講義では、方程式の根の可換性や置換を主体とする古典的なGalois理論の記述スタイルから離れ、多項式環のイデアル論、テンソル積の完全分解（分裂）、および自己準同型環の構造定理（Jacobson-Bourbakiの定理）を用いた、現代的な抽象代数学におけるGalois理論の再構築を行う。
    
    数式や証明のステップを極めて丁寧に記述し、基本概念の定義からGaloisの基本定理（Galois対応の一対一性）の完成までを自己完結的 (self-contained) に解説する。
\end{frame}

\section{基本概念の定義と準備}
\begin{frame}[allowframebreaks]{1. 基本概念の定義と準備}
    本講義において、集合の包含関係は $\subset$ を用い、真の包含関係である必要はないものとする。空集合は $\varnothing$ で表し、集合 $A$ から集合 $B$ を除く差集合は $A \smallsetminus B$ と表記する。

    \begin{block}{定義 1.1 (有限次Galois拡大とGalois群)}
        体 $K$ の有限次代数拡大 $L$ が正規拡大かつ分離拡大であるとき、$L/K$ を有限次Galois拡大 (finite Galois extension) と呼ぶ。また、$K$ の各元を固定する $L$ の自己同型全体のなす群をGalois群 (Galois group) と呼び、$G = \mathrm{Gal}(L/K)$ で表す。
    \end{block}

    \begin{block}{定義 1.2 (テンソル積代数)}
        $K$ を体とし、$A, B$ を $K$ 上の可換代数とする。$A$ と $B$ の $K$ 上のベクトル空間としてのテンソル積 (tensor product) $A \otimes_K B$ に対し、乗法を
        \[ (a_1 \otimes b_1)(a_2 \otimes b_2) = (a_1 a_2) \otimes (b_1 b_2) \quad (a_1, a_2 \in A, \; b_1, b_2 \in B) \]
        により定めた環構造を $K$ 上のテンソル積代数と呼ぶ。
    \end{block}

    \begin{block}{定義 1.3 (ひねり群環 / 接合積)}
        有限次Galois拡大 $L/K$ とそのGalois群 $G = \mathrm{Gal}(L/K)$ に対し、各 $\sigma \in G$ に対応する記号を $\sigma$ そのもので表す。$L$ 上のベクトル空間としての直和
        \[ L[G] = \bigoplus_{\sigma \in G} L\sigma \]
        に対し、乗法規則を $(a\sigma)(b\tau) = a\sigma(b)\sigma\tau \; (a, b \in L, \; \sigma, \tau \in G)$ および分配法則によって定義した非可換環をひねり群環 (skew group ring) または接合積 (crossed product) と呼ぶ。
    \end{block}

    \begin{block}{定義 1.4 (自己準同型環)}
        $L$ を $K$ 上のベクトル空間とみなしたとき、$L$ から $L$ への $K$-線形写像全体のなす集合を $\mathrm{End}_K(L)$ と表す。写像の和と合成（積）に関して $\mathrm{End}_K(L)$ は非可換環となり、これを $K$ 上の自己準同型環 (endomorphism ring) と呼ぶ。
    \end{block}
\end{frame}

\section{予備的補題とその証明}
\begin{frame}[allowframebreaks]{2. 予備的補題とその証明}
    \begin{block}{補題 2.1 (原始元定理)}
        $L/K$ が有限次分離拡大であるならば、ある単一の元 $\alpha \in L$ が存在して $L = K(\alpha)$ と表される。
    \end{block}
    \begin{block}{補題 2.1の証明}
        $K$ が有限体の場合、拡大体 $L$ も有限体。有限体の乗法群は巡回群であるため、生成元 $g$ をとれば $L = K(g)$。\\
        $K$ が無限体の場合、$L = K(\beta, \gamma)$ の場合を示せば帰納法により一般に帰着される。$\beta, \gamma$ の最小多項式を $p(X), q(X)$ とし、根をそれぞれ $\beta=\beta_1, \dots, \beta_r$ および $\gamma=\gamma_1, \dots, \gamma_s$ とする。拡大の分離性より相異なる。
        任意の $i$ および $j \geq 2$ に対して $\beta_i + c\gamma_j \neq \beta_1 + c\gamma_1$ となる $c \in K$ を選ぶ（有限個の禁止点を除く）。$\alpha = \beta + c\gamma \in L$ と置く。\\
        多項式 $h(X) = p(\alpha - cX)$ は $h(\gamma) = p(\beta) = 0$。$h(X)$ と $q(X)$ の共通根を調べると、$c$ の選び方により共通根は $\gamma$ のみとなる。よって $\gamma \in K(\alpha)$、$\beta = \alpha - c\gamma \in K(\alpha)$ となり、$L = K(\alpha)$ である。
    \end{block}

    \begin{block}{補題 2.2 (中国剰余定理)}
        可換環 $R$ のイデアル $I_1, I_2, \dots, I_m$ がどの2つも互いに素であるとする。このとき、以下の環同型が成り立つ。
        \[ R/\left( \bigcap_{k=1}^m I_k \right) \cong \prod_{k=1}^m R/I_k \]
    \end{block}
    \begin{block}{補題 2.2の証明}
        自然な環準同型写像 $\phi: R \to \prod_{k=1}^m R/I_k \; (x \mapsto (x \bmod I_1, \dots))$ を考える。核は明らかに $\bigcap_{k=1}^m I_k$。\\
        $m=2$ のとき、$e_1 + e_2 = 1 \; (e_1 \in I_1, e_2 \in I_2)$ より、$x = r_2 e_1 + r_1 e_2$ と置けば $\phi(x) = (r_1 \bmod I_1, r_2 \bmod I_2)$ となり全射。\\
        $m > 2$ のとき、各 $k \geq 2$ に対して $u_k + v_k = 1 \; (u_k \in I_1, v_k \in I_k)$ が存在する。このとき
        \[ 1 = \prod_{k=2}^m (u_k + v_k) \in I_1 + \prod_{k=2}^m I_k \subset I_1 + \bigcap_{k=2}^m I_k \]
        より $I_1$ と $\bigcap_{k=2}^m I_k$ は互いに素。帰納法により全射性が示される。
    \end{block}

    \begin{block}{補題 2.3 (Dedekindの指標の独立性定理)}
        相異なる体準同型 $\sigma_1, \dots, \sigma_m: L \to M$ は、$M$ 上線形独立である。すなわち、$\sum_{i=1}^m c_i \sigma_i(x) = 0 \; (\forall x \in L)$ を満たす $c_i \in M$ はすべて $0$。
    \end{block}
    \begin{block}{補題 2.3の証明}
        写像の個数 $m$ による帰納法。$m=1$ のとき自明。\\
        $m-1$ で成立と仮定。$\sum_{i=1}^m c_i \sigma_i(x) = 0$ とし、$c_i$ に非零が存在するとする。ある $y \in L$ で $\sigma_1(y) \neq \sigma_m(y)$。自変数を $yx$ に代入した式から、元の式に $\sigma_m(y)$ を乗じたものを引くと、
        \[ \sum_{i=1}^{m-1} c_i (\sigma_i(y) - \sigma_m(y))\sigma_i(x) = 0 \]
        帰納法の仮定より各係数は $0$。特に $c_1(\sigma_1(y) - \sigma_m(y)) = 0$。$c_1 \neq 0$ と仮定していたため $\sigma_1(y) = \sigma_m(y)$ となり矛盾。
    \end{block}
\end{frame}

\section{テンソル積の完全分解}
\begin{frame}[allowframebreaks]{3. 主定理1：テンソル積 $L \otimes_K L$ の完全分解}
    \begin{block}{定理 3.1}
        $L/K$ を有限次Galois拡大とし、$G = \mathrm{Gal}(L/K)$ とする。このとき、以下の写像 $\Psi$ は可換環の同型写像である。
        \[ \Psi: L \otimes_K L \longrightarrow \prod_{\sigma \in G} L, \quad a \otimes b \longmapsto (a\sigma(b))_{\sigma \in G} \]
    \end{block}
    \begin{block}{定理 3.1の証明}
        原始元定理より $L = K(\alpha)$。$\alpha$ の最小多項式を $f(X) \in K[X]$ とすると $L \cong K[X]/(f(X))$。Galois拡大であるから $f(X)$ は完全に分解され $f(X) = \prod_{\sigma \in G} (X - \sigma(\alpha))$。\\
        テンソル積の性質より同型 $\theta$ を得る。
        \[ \theta: L \otimes_K L \cong L[X]/(f(X)), \quad a \otimes g(\alpha) \mapsto a \cdot g(X) \pmod{f(X)} \]
        中国剰余定理より同型 $\omega$ を得る。
        \[ \omega: L[X]/(f(X)) \cong \prod_{\sigma \in G} L[X]/(X - \sigma(\alpha)) \cong \prod_{\sigma \in G} L \]
        合成写像 $\omega \circ \theta$ の行き先を計算すると、$\omega(a \cdot g(X)) = (a \cdot g(\sigma(\alpha)))_{\sigma \in G} = (a\sigma(b))_{\sigma \in G}$ となり、これは $\Psi$ に一致する。よって同型。
    \end{block}
\end{frame}

\section{応用としての次数の決定}
\begin{frame}[allowframebreaks]{4. 応用としての次数の決定}
    \begin{block}{定理 4.1}
        有限次Galois拡大 $L/K$ に対し、拡大次数 $[L:K]$ はGalois群の位数 $|G|$ と一致する。
    \end{block}
    \begin{block}{定理 4.1の証明}
        同型 $\Psi: L \otimes_K L \xrightarrow{\cong} \prod_{\sigma \in G} L$ の両辺に、左から $L$ の元 $x$ を掛ける作用を定義し、左 $L$-ベクトル空間とみなす。
        $\Psi$ はこの作用と可換（$\Psi(x \cdot (a \otimes b)) = x \cdot \Psi(a \otimes b)$）であるため、左 $L$-ベクトル空間としての同型。よって $L$ 上の次元が一致する。\\
        左辺について、$K$ 上の基底 $v_1, \dots, v_n \; (n = [L:K])$ を用いると、$1 \otimes v_i$ は $L$ 上の基底となるため、$\dim_L(L \otimes_K L) = [L:K]$。\\
        右辺について、$|G|$ 個の $L$ の直積であるため次元はそのまま成分の個数、$\dim_L(\prod_{\sigma \in G} L) = |G|$。次元の比較より $[L:K] = |G|$。
    \end{block}
\end{frame}

\section{自己準同型環の同型}
\begin{frame}[allowframebreaks]{5. 主定理2：自己準同型環の同型}
    \begin{block}{定理 5.1}
        有限次Galois拡大 $L/K$ に対し、ひねり群環 $L[G]$ から、自己準同型環 $\mathrm{End}_K(L)$ への以下の写像 $\Phi$ は、環としての同型写像である。
        \[ \Phi: L[G] \longrightarrow \mathrm{End}_K(L), \quad \Phi\left( \sum_{\sigma \in G} a_\sigma \sigma \right)(x) = \sum_{\sigma \in G} a_\sigma \sigma(x) \]
    \end{block}
    \begin{block}{定理 5.1の証明}
        積の保存は $\Phi(a\sigma(b)\sigma\tau)(x) = a\sigma(b\tau(x))$ より従う。和と分配法則も自明であり、$\Phi$ は環準同型。\\
        $\Phi(\sum a_\sigma \sigma) = 0$ とすると全 $x$ で $\sum a_\sigma \sigma(x) = 0$。Dedekindの指標の独立性（補題 2.3）より全 $\sigma$ で $a_\sigma = 0$ となり、単射。\\
        次元を比較する。定理 4.1 より $|G| = n = [L:K]$ と置くと、$\dim_K(L[G]) = n \times n = n^2$。また $\dim_K(\mathrm{End}_K(L)) = n^2$。次元の一致と単射性より全射。よって同型。
    \end{block}
\end{frame}

\section{中心化環と包含関係の反転機構}
\begin{frame}[allowframebreaks]{6. 中心化環と包含関係の反転機構}
    \begin{block}{定義 6.1 (中心化環)}
        環 $A$ と部分集合 $X$ に対し、$X$ の全元と可換な $A$ の元全体のなす部分環 $C_A(X) = \{ a \in A \mid ax = xa \quad (\forall x \in X) \}$ を中心化環と呼ぶ。
    \end{block}

    $A = \mathrm{End}_K(L)$ を考える。$L$ の元 $a$ による左乗算作用を $\lambda_a(x) = ax$ とし、これを $L_L \subset A$ と表す。中間体 $M$ に対しても $M_L \subset A$ が定まる。

    \begin{block}{補題 6.2}
        中間体 $M$ に対し、$C_A(M_L) = \mathrm{End}_M(L)$ が成り立つ。
    \end{block}
    \begin{block}{補題 6.2の証明}
        $f \in C_A(M_L)$ の条件は、任意の $m \in M$ に対し $f \circ \lambda_m = \lambda_m \circ f$。これは $f(mx) = mf(x)$ と同値であり、$f$ が $M$-線形写像である定義そのもの。
    \end{block}
\end{frame}

\section{Galoisの基本定理の現代的証明}
\begin{frame}[allowframebreaks]{7. Galoisの基本定理の現代的証明}
    \begin{block}{定理 7.1 (Galoisの基本定理)}
        中間体全体の集合 $\mathcal{M}$ と部分群全体の集合 $\mathcal{H}$ の間に、互いに逆の一対一対応が存在する。
        \[ \mathcal{M} \longrightarrow \mathcal{H} \; (M \mapsto \mathrm{Gal}(L/M)), \quad \mathcal{H} \longrightarrow \mathcal{M} \; (H \mapsto L^H) \]
    \end{block}
    \begin{block}{定理 7.1の証明}
        同型 $\Phi: L[G] \cong A = \mathrm{End}_K(L)$ のもと、部分環 $\mathrm{End}_M(L)$ に対応する $L[G]$ の元 $\sum a_\sigma \sigma$ を調べる。
        $\sum a_\sigma \sigma(mx) = m \sum a_\sigma \sigma(x)$ より $a_\sigma \neq 0 \Rightarrow \sigma \in \mathrm{Gal}(L/M)$。
        よって $\mathrm{End}_M(L) \cong L[\mathrm{Gal}(L/M)]$。\\
        逆に、部分群 $H$ に対し $L[H]$ の $L_L$ における中心化環を考えると、$x\sigma(y) = \sigma(xy)$ より $x = \sigma(x)$ となり、$C_{L_L}(L[H]) = (L^H)_L$。\\
        (1) $M$ から出発して $H = \mathrm{Gal}(L/M)$ とすると、二重中心化環の一般性質から $C_A(\mathrm{End}_M(L)) = M_L$。これと $L_L$ の交叉から $L^{\mathrm{Gal}(L/M)} = M$。\\
        (2) $H$ から出発して $M = L^H$ とすると、$L[H] \subset L[\mathrm{Gal}(L/M)]$。次元を計算すると双方ともに等しくなるため、位数が一致し $H = \mathrm{Gal}(L/L^H)$ が従う。
    \end{block}
\end{frame}

\section{現代数学における視点}
\begin{frame}[allowframebreaks]{8. 現代数学における視点}
    本講義で展開した $L \otimes_K L \cong \prod L$ および $L[G] \cong \mathrm{End}_K(L)$ を経由する証明手法は、方程式の個別の根の計算を完全にバイパスし、空間の持つ対称性そのものを「環の分解」および「自己準同型環の表現」という強力な道具で記述するものである。

    このアプローチは、N. Jacobsonらによって \textbf{Jacobson-Bourbakiの定理} として一般化され、斜体拡大のガロア理論の構築へと繋がった。さらに、A. Grothendieckによる \textbf{エタール被覆 (étale coverings)} の理論の基礎となり、現代代数幾何学におけるエタール基本群の定義へと直接的に直結している。
\end{frame}

\end{document}